% GENERATED FILE. Edit canonical lessons, then run npm run generate:books.
% canonical-source-hash: fnv1a64:3e78ba16659cc824
% canonical-lessons: GE-C424-gasthof, GE-C424-ferienwohnung, GE-C424-uebernachten, GE-C424-uebernachtung, GE-C424-halbpension

\chapter{Wo schlafen wir}
\label{ch:ge-wo-schlafen-wir}
\hlchaptermodality{424}

\begin{chapteropening}
\textbf{By the end of this chapter:} \emph{I can talk about places to stay: the inn, the holiday flat, staying the night, the overnight stay and half board.}

Complete the last lesson of chapter 424: I can talk about places to stay: the inn, the holiday flat, staying the night, the overnight stay and half board.
\end{chapteropening}

\section[der Gasthof]{der Gasthof — the inn, the country hotel}

\label{lesson:GE-C424-gasthof}



\begin{quote}
Before the new one: say the German for the sailor, then the German for to capsize.
\end{quote}

\subsection*{You'll want to know: der Gasthof}
\textbf{der Gasthof} — \textquotedblleft{}the inn, the country hotel\textquotedblright{}.

\begin{quote}
Im Dorf gibt es einen alten Gasthof mit gutem Essen.
\end{quote}

\begin{quote}
\emph{In the village there is an old inn with good food.}
\end{quote}

Plural: \emph{die Gasthöfe}. Common in the south and in Austria.

\begin{cousinweb}
\emph{Gasthof} is built from \emph{Gast} 'guest', inherited from Proto-Germanic \emph{gastiz} and a cousin of English \emph{guest} (EtymDB), and \emph{Hof} 'yard, farm, court', also an inherited word.
\end{cousinweb}

\begin{grammarlens}[title={What you've built}]
One more word for telling, explaining and describing at B1.
\end{grammarlens}

\subsection*{Your turn}
\begin{itemize}
\raggedright
  \item \emph{Say it:} \emph{der Gasthof}
  \item \emph{Say it:} \emph{der Gasthof}, once more
  \item \emph{Say it:} a sentence of your own with \emph{der Gasthof}
  \item \emph{Recall:} say the German for the sailor, then the German for to capsize, then say \emph{der Gasthof} again
\end{itemize}

\begin{tcolorbox}[breakable,colback=teal!4,colframe=teal!35!black,title={Before you move on}]
What does der Gasthof mean? (\textquotedblleft{}The inn, the country hotel\textquotedblright{}.) Say it once more.
\end{tcolorbox}
\section[die Ferienwohnung]{die Ferienwohnung — the holiday flat, the holiday apartment}

\label{lesson:GE-C424-ferienwohnung}



\begin{quote}
Before the new one: say the German for to capsize, then the German for the inn.
\end{quote}

\subsection*{You'll want to know: die Ferienwohnung}
\textbf{die Ferienwohnung} — \textquotedblleft{}the holiday flat, the holiday apartment\textquotedblright{}.

\begin{quote}
Wir mieten eine Ferienwohnung am See.
\end{quote}

\begin{quote}
\emph{We are renting a holiday flat by the lake.}
\end{quote}

Plural: \emph{die Ferienwohnungen}.

\begin{cousinweb}
\emph{Ferienwohnung} is built from \emph{Ferien} 'holidays', a learned borrowing from Latin \emph{fēriae} 'days of rest, festivals' (Lewis \& Short), and \emph{Wohnung} 'flat', from \emph{wohnen} 'to live', inherited from Old High German \emph{wonēn} (EtymDB).
\end{cousinweb}

\begin{grammarlens}[title={What you've built}]
One more word for telling, explaining and describing at B1.
\end{grammarlens}

\subsection*{Your turn}
\begin{itemize}
\raggedright
  \item \emph{Say it:} \emph{die Ferienwohnung}
  \item \emph{Say it:} \emph{die Ferienwohnung}, once more
  \item \emph{Say it:} a sentence of your own with \emph{die Ferienwohnung}
  \item \emph{Recall:} say the German for to capsize, then the German for the inn, then say \emph{die Ferienwohnung} again
\end{itemize}

\begin{tcolorbox}[breakable,colback=teal!4,colframe=teal!35!black,title={Before you move on}]
What does die Ferienwohnung mean? (\textquotedblleft{}The holiday flat, the holiday apartment\textquotedblright{}.) Say it once more.
\end{tcolorbox}
\section[übernachten]{übernachten — to stay the night, to spend the night}

\label{lesson:GE-C424-uebernachten}



\begin{quote}
Before the new one: say the German for the inn, then the German for the holiday flat.
\end{quote}

\subsection*{You'll want to know: übernachten}
\textbf{übernachten} — \textquotedblleft{}to stay the night, to spend the night\textquotedblright{}.

\begin{quote}
Auf dem Weg in den Süden übernachten wir in einem kleinen Hotel.
\end{quote}

\begin{quote}
\emph{On the way south we spend the night in a small hotel.}
\end{quote}

Inseparable: \emph{wir übernachten}, \emph{wir haben übernachtet}.

\begin{cousinweb}
\emph{übernachten} is built from \emph{über} 'over' and \emph{Nacht} 'night', both inherited Germanic words: \emph{Nacht} goes back to Old High German \emph{naht} (EtymDB).
\end{cousinweb}

\begin{grammarlens}[title={What you've built}]
One more word for telling, explaining and describing at B1.
\end{grammarlens}

\subsection*{Your turn}
\begin{itemize}
\raggedright
  \item \emph{Say it:} \emph{übernachten}
  \item \emph{Say it:} \emph{übernachten}, once more
  \item \emph{Say it:} a sentence of your own with \emph{übernachten}
  \item \emph{Recall:} say the German for the inn, then the German for the holiday flat, then say \emph{übernachten} again
\end{itemize}

\begin{tcolorbox}[breakable,colback=teal!4,colframe=teal!35!black,title={Before you move on}]
What does übernachten mean? (\textquotedblleft{}To stay the night, to spend the night\textquotedblright{}.) Say it once more.
\end{tcolorbox}
\section[die Übernachtung]{die Übernachtung — the overnight stay, the night (in a hotel)}

\label{lesson:GE-C424-uebernachtung}



\begin{quote}
Before the new one: say the German for the holiday flat, then the German for to stay the night.
\end{quote}

\subsection*{You'll want to know: die Übernachtung}
\textbf{die Übernachtung} — \textquotedblleft{}the overnight stay, the night (in a hotel)\textquotedblright{}.

\begin{quote}
Eine Übernachtung mit Frühstück kostet achtzig Euro.
\end{quote}

\begin{quote}
\emph{One night with breakfast costs eighty euros.}
\end{quote}

Plural: \emph{die Übernachtungen}.

\begin{cousinweb}
\emph{Übernachtung} is built on \emph{übernachten} 'to stay the night' with the ending \emph{-ung}; the verb comes from \emph{über} and \emph{Nacht} 'night'.
\end{cousinweb}

\begin{grammarlens}[title={What you've built}]
One more word for telling, explaining and describing at B1.
\end{grammarlens}

\subsection*{Your turn}
\begin{itemize}
\raggedright
  \item \emph{Say it:} \emph{die Übernachtung}
  \item \emph{Say it:} \emph{die Übernachtung}, once more
  \item \emph{Say it:} a sentence of your own with \emph{die Übernachtung}
  \item \emph{Recall:} say the German for the holiday flat, then the German for to stay the night, then say \emph{die Übernachtung} again
\end{itemize}

\begin{tcolorbox}[breakable,colback=teal!4,colframe=teal!35!black,title={Before you move on}]
What does die Übernachtung mean? (\textquotedblleft{}The overnight stay, the night (in a hotel)\textquotedblright{}.) Say it once more.
\end{tcolorbox}
\section[die Halbpension]{die Halbpension — the half board}

\label{lesson:GE-C424-halbpension}



\begin{quote}
Before the new one: say the German for to stay the night, then the German for the overnight stay.
\end{quote}

\subsection*{You'll want to know: die Halbpension}
\textbf{die Halbpension} — \textquotedblleft{}the half board\textquotedblright{}.

\begin{quote}
Mit Halbpension bekommen wir Frühstück und Abendessen.
\end{quote}

\begin{quote}
\emph{With half board we get breakfast and dinner.}
\end{quote}

No plural. \emph{Die Vollpension} — full board.

\begin{cousinweb}
\emph{Halbpension} is built from \emph{halb} 'half', inherited from Proto-Germanic \emph{halbaz} (EtymDB), and \emph{Pension} 'board and lodging', borrowed from French \emph{pension} (EtymDB), which goes back to Latin \emph{pēnsiō} 'payment' (Lewis \& Short).
\end{cousinweb}

\begin{grammarlens}[title={What you've built}]
One more word for telling, explaining and describing at B1.
\end{grammarlens}

\subsection*{Your turn}
\begin{itemize}
\raggedright
  \item \emph{Say it:} \emph{die Halbpension}
  \item \emph{Say it:} \emph{die Halbpension}, once more
  \item \emph{Say it:} a sentence of your own with \emph{die Halbpension}
  \item \emph{Recall:} say the German for to stay the night, then the German for the overnight stay, then say \emph{die Halbpension} again
\end{itemize}

\begin{tcolorbox}[breakable,colback=teal!4,colframe=teal!35!black,title={Before you move on}]
What does die Halbpension mean? (\textquotedblleft{}The half board\textquotedblright{}.) Say it once more.
\end{tcolorbox}
